\begin{table}[htbp]
\centering
\caption{Efecto del número de clases sobre la efectividad de la aumentación (pond. suave vs SMOTE)}
\label{tab:nclasses_pattern}
\begin{tabular}{rlrrrll}
\toprule
\textbf{$N$ clases} & \textbf{Conjunto(s)} & \textbf{$\Delta$ (pp)} & \textbf{Std} & \textbf{Vic.\%} & \textbf{$N$ pares} \\
\midrule
2   & sms\_spam                      & +3.18 & 3.24 & 93.3\%  & 15 \\
3   & hate\_speech\_davidson          & +1.86 & 2.97 & 66.7\%  & 15 \\
4   & 20newsgroups, ag\_news          & +1.70 & 3.02 & 73.3\%  & 30 \\
\rowcolor{green!20}
6   & emotion, trec6                 & \textbf{+5.42} & 3.63 & \textbf{100.0\%} & 24 \\
14  & dbpedia14                      & +0.40 & 0.61 & 80.0\%  & 15 \\
20  & 20newsgroups\_20class           & +2.36 & 1.55 & 100.0\% & 15 \\
77  & banking77                      & $-$0.99 & 0.52 & 0.0\% & 9 \\
150 & clinc150                       & +0.01 & 0.53 & 44.4\%  & 9 \\
\midrule
\multicolumn{6}{l}{\small Correlación de Spearman: $r=-0.260$, $p=0.003$. Más clases $\Rightarrow$ menores ganancias.} \\
\multicolumn{6}{l}{\small Rango óptimo: 2--20 clases. Las ganancias disminuyen a partir de 77 clases.} \\
\bottomrule
\end{tabular}
\end{table}
